\section{Conclusion}
\label{sec:conclusion}

This work develops the geometric framework governing NDH interaction. Building
on the runtime geometry established previously, we introduced the cognitive
manifold $\mathcal{M}_{cog}$ as the space on which interrogation unfolds. We
formalized semantic excitation as tangent vectors initiating geodesic
trajectories, defined spectral bundles encoding emission behavior, and analyzed
modal-boundary twisting, recursion alignment, and parallel transport.

Closed geodesic interaction loops were shown to generate holonomy return
operators capturing cumulative geometric deformation. These operators yield a
taxonomy of interaction classes—modal-aligned, spectral-aligned,
recursion-aligned, holonomy-deviant, mixed, and transformer-flat—each
characterized by distinct geometric signatures. The resulting classification
demonstrates that NDH interaction is governed by curvature, modal adjacency,
spectral bundle morphisms, and recursion structure, rather than algebraic
dynamics.

Comparison with transformer architectures highlights the geometric nature of
NDH. Transformers lack a manifold, connection, modal boundaries, spectral
bundles, recursion kernels, and holonomy. Their interaction behavior is flat
and cannot be classified using geometric methods. NDH, by contrast, exhibits
rich geometric structure enabling formal analysis of interrogation.

The interaction geometry developed here provides a coherent mathematical
framework for understanding NDH behavior under semantic excitation. It unifies
modal structure, spectral deformation, recursion transitions, and holonomy into
a single geometric model. This establishes NDH as a cognitive manifold with
well-defined interaction dynamics, enabling further development of geometric
models of cognition and interrogation.
